\documentclass[12pt,a4paper]{article}
\usepackage[utf8]{inputenc}
\usepackage[T1]{fontenc}
\usepackage[french]{babel}
\usepackage{geometry}
\usepackage{xcolor}
\usepackage{fancyhdr}
\usepackage{array}

\geometry{margin=2cm}
\pagestyle{fancy}
\fancyhf{}
\fancyhead[L]{CEJM - Fiches de Révision}
\fancyhead[R]{Vue d'ensemble}
\fancyfoot[C]{\thepage}
\setlength{\headheight}{15pt}

% Couleurs
\definecolor{bleu}{RGB}{0,90,170}
\definecolor{vert}{RGB}{0,150,0}
\definecolor{rouge}{RGB}{200,0,0}
\definecolor{gris}{RGB}{240,240,240}

\title{\textbf{CEJM - Fiches de Révision}\\[2mm]\large Vue d'ensemble des 4 chapitres}
\author{}
\date{}

\begin{document}
\maketitle

\section*{Structure du programme CEJM}

\subsection*{Thème 1 : L'intégration de l'entreprise dans son environnement}

\begin{center}
\begin{tabular}{|p{3cm}|p{6cm}|p{6cm}|}
\hline
\textbf{Chapitre} & \textbf{Problématique} & \textbf{Concepts clés} \\
\hline
\textbf{Chapitre 1} & Comment s'établissent les relations entre l'entreprise et son environnement économique ? & Agents économiques, marchés, flux réels/monétaires, circuit économique \\
\hline
\textbf{Chapitre 2} & Comment fonctionnent les mécanismes de marché et comment se forme le prix d'équilibre ? & Offre/demande, prix d'équilibre, élasticité, concurrence \\
\hline
\textbf{Chapitre 3} & Comment se forme juridiquement un contrat et quelles sont les conditions de sa validité ? & Formation du contrat, conditions de validité, vices du consentement \\
\hline
\textbf{Chapitre 4} & Comment créer une entreprise et mesurer sa performance ? & Création d'entreprise, formes juridiques, performance \\
\hline
\end{tabular}
\end{center}

\section{Chapitre 1 : Les agents économiques}

\subsection{L'essentiel}
\begin{itemize}
    \item \textbf{4 agents principaux} : Ménages, Entreprises, Banques, Administrations publiques
    \item \textbf{3 marchés fondamentaux} : Biens/services, Travail, Capitaux
    \item \textbf{Circuit économique} : flux réels <-> flux monétaires
    \item \textbf{Interdépendance} : tous les agents sont liés par des échanges
\end{itemize}

\subsection{Mots-clés indispensables}
Agents économiques, Production marchande/non marchande, Circuit économique, Marché des capitaux

\section{Chapitre 2 : Le fonctionnement des marchés}

\subsection{L'essentiel}
\begin{itemize}
    \item \textbf{Loi de l'offre} : prix augmente -> offre augmente
    \item \textbf{Loi de la demande} : prix augmente -> demande diminue
    \item \textbf{Équilibre} : intersection offre/demande
    \item \textbf{Autorégulation} : ajustement automatique par les prix
\end{itemize}

\subsection{Mots-clés indispensables}
Offre, Demande, Prix d'équilibre, Élasticité, Concurrence parfaite, Monopole

\section{Chapitre 3 : La formation du contrat}

\subsection{L'essentiel}
\begin{itemize}
    \item \textbf{4 conditions de validité} : Consentement, Capacité, Objet, Cause
    \item \textbf{Formation} : Offre + Acceptation = Contrat
    \item \textbf{Vices du consentement} : Erreur, Dol, Violence
    \item \textbf{Sanctions} : Nullité absolue ou relative
\end{itemize}

\subsection{Mots-clés indispensables}
Contrat, Offre, Acceptation, Consentement, Vices du consentement, Nullité

\section{Chapitre 4 : La création d'entreprise}

\subsection{L'essentiel}
\begin{itemize}
    \item \textbf{Étapes de création} : Idée -> Business plan -> Financement -> Forme juridique
    \item \textbf{Formes principales} : Entreprise individuelle, SARL, SAS, SA
    \item \textbf{Performance} : Commerciale, Financière, Sociale
    \item \textbf{Indicateurs clés} : CA, Résultat, Rentabilité, Parts de marché
\end{itemize}

\subsection{Mots-clés indispensables}
Entrepreneur, Business plan, SARL/SAS, Chiffre d'affaires, Rentabilité, Performance

\section{Liens entre les chapitres}

\subsection{Cohérence du programme}
\begin{itemize}
    \item \textbf{Ch1 -> Ch2} : Les agents économiques interagissent sur les marchés
    \item \textbf{Ch2 -> Ch3} : Les échanges sur les marchés nécessitent des contrats
    \item \textbf{Ch3 -> Ch4} : Les contrats permettent la création et le développement des entreprises
    \item \textbf{Ch4 -> Ch1} : Les entreprises créées deviennent des agents économiques
\end{itemize}

\section{Méthodes de révision}

\subsection{Pour les définitions}
\begin{itemize}
    \item Apprendre par cœur les définitions clés
    \item Faire des fiches bristol avec terme/définition
    \item S'entraîner à donner des exemples concrets
\end{itemize}

\subsection{Pour les mécanismes}
\begin{itemize}
    \item Dessiner les schémas (circuit économique, courbes offre/demande)
    \item Expliquer les enchaînements cause/effet
    \item Utiliser des exemples d'actualité économique
\end{itemize}

\subsection{Pour les applications}
\begin{itemize}
    \item Analyser des cas d'entreprises réelles
    \item Résoudre des exercices de calcul (performance)
    \item Rédiger des synthèses structurées
\end{itemize}

\section{Conseils pour l'examen}

\subsection{Structure type d'une réponse}
\begin{enumerate}
    \item \textbf{Introduction} : Définir les termes, annoncer le plan
    \item \textbf{Développement} : 2-3 parties avec exemples
    \item \textbf{Conclusion} : Synthèse et ouverture
\end{enumerate}

\subsection{Erreurs à éviter}
\begin{itemize}
    \item Confondre les notions (offre/demande, CA/résultat)
    \item Oublier les exemples concrets
    \item Ne pas relier les concepts entre eux
    \item Négliger l'actualité économique
\end{itemize}

\subsection{Points à valoriser}
\begin{itemize}
    \item Maîtrise du vocabulaire technique
    \item Exemples d'entreprises actuelles
    \item Liens avec l'actualité économique
    \item Schémas et graphiques clairs
\end{itemize}

\section{Sujets transversaux importants}

\subsection{L'entreprise au cœur du système}
L'entreprise est présente dans tous les chapitres comme acteur central de l'économie.

\subsection{Les relations contractuelles}
Tous les échanges économiques se matérialisent par des contrats (travail, vente, financement).

\subsection{La performance et la concurrence}
Les entreprises évoluent sur des marchés concurrentiels et doivent optimiser leur performance.

\subsection{Le rôle de l'État}
L'État intervient comme agent économique, régulateur des marchés et cadre juridique des contrats.

\end{document}